% Generated draft: requires mathematical and source review.
\documentclass[11pt]{article}
\usepackage[margin=25mm]{geometry}
\usepackage{fontspec}
\setmainfont{Latin Modern Roman}
\usepackage{amsmath,amssymb,mathtools,unicode-math}
\setmathfont{Latin Modern Math}
\usepackage{enumitem,needspace,xurl,hyperref}
\setlength{\parindent}{0pt}
\setlength{\parskip}{5pt}
\setlength{\emergencystretch}{3em}
\newcommand{\N}{\mathbb{N}}
\newcommand{\Z}{\mathbb{Z}}
\newcommand{\Q}{\mathbb{Q}}
\newcommand{\R}{\mathbb{R}}
\newcommand{\C}{\mathbb{C}}
\newcommand{\F}{\mathbb{F}}
\newcommand{\A}{\mathbb{A}}
\newcommand{\PP}{\mathbb{P}}
\newcommand{\E}{\mathbb{E}}
\newcommand{\eps}{\varepsilon}
\newcommand{\dd}{\,\mathrm{d}}
\begin{document}
\section*{UnsolvedMath ID: 701275}
\textbf{Source problem code:} catalogue entry
\subsection{Definitions and mathematical statement}
For an integer \(n\geq 4\), let \(\mathfrak B_n\) be the braid group on \(n\) strands, with presentation
\[
\mathfrak B_n=\left\langle \sigma_1,\ldots,\sigma_{n-1}\ \middle|\ \sigma_i\sigma_{i+1}\sigma_i=\sigma_{i+1}\sigma_i\sigma_{i+1}\ (1\leq i\leq n-2),\quad \sigma_i\sigma_j=\sigma_j\sigma_i\ (|i-j|\geq2)\right\rangle.
\]
Under the usual finite-branching interpretation of a regular rooted tree, choose an integer \(d\geq2\) and an alphabet \(X\) with \(|X|=d\). The tree \(X^*\) has all finite words over \(X\) as vertices, the empty word as root, and edges joining \(v\) to \(vx\) for \(x\in X\). Its automorphisms preserve the root and adjacency. For \(g\in\operatorname{Aut}(X^*)\) and \(v\in X^*\), the section \(g|_v\in\operatorname{Aut}(X^*)\) is determined by
\[
 g(vw)=g(v)\,g|_v(w)\qquad(w\in X^*).
\]
An automorphism \(g\) is finite-state if \(\{g|_v:v\in X^*\}\) is finite. An action \(\rho:\mathfrak B_n\to\operatorname{Aut}(X^*)\) is faithful if \(\ker\rho=\{1\}\), and is by finite-state automorphisms if every \(\rho(b)\), for \(b\in\mathfrak B_n\), is finite-state.

The question is whether, for every integer \(n\geq4\), there exist an integer \(d\geq2\), an alphabet \(X\) of cardinality \(d\), and an injective homomorphism
\[
 \rho_n:\mathfrak B_n\longrightarrow\operatorname{Aut}(X^*)
\]
such that
\[
 \forall b\in\mathfrak B_n,\qquad
 \left|\{\rho_n(b)|_v:v\in X^*\}\right|<\infty.
\]
Here the tree degree and action may depend on \(n\); no common bound on the number of states is imposed. The source does not explicitly specify the tree degree or finite branching, so the finite-alphabet convention used above requires confirmation against its cited definitions. The source also states that such an action is known for \(\mathfrak B_3\); this is preserved as supplied background, not as an independently verified result.
\subsection{Short English statement}
For each \(n\geq4\), can the braid group on \(n\) strands act faithfully on a regular rooted tree so that each group element acts with only finitely many distinct sections? The formulation uses the conventional finite-branching interpretation; the source leaves the tree degree unspecified. It reports that the corresponding action exists for \(\mathfrak B_3\).
\subsection{Sources}
\begin{itemize}
\item Supplied local UnsolvedMath registry record, permanent ID \(701275\), exact source code \texttt{catalogue entry}, title \emph{Kourovka Notebook Problem 16.84}. The registry labels the problem \emph{open}; that label is supplied metadata, not an independently verified literature assessment. The record gives proposal year 2006 and upstream revision \texttt{d371ee917ac9647fbb06cd6fc529deec6f3dda82}.
\item \emph{The Kourovka Notebook (Unsolved Problems in Group Theory)}, Problem 16.84, 16th issue (2006); the supplied background also identifies the 21st issue edition (2026), \texttt{arXiv:1401.0300v48}. The collection description names E. I. Khukhro and V. D. Mazurov as editors. These publication details are supplied attribution and have not been independently checked.
\item The registry attributes the problem to V. I. Sushchanskiĭ, but its proposer field contains a truncated TeX rendering of the surname. The complete surname rendering used here occurs in the reference supplied with the statement.
\item Reference supplied for definitions: R. I. Grigorchuk, V. V. Nekrashevich, V. I. Sushchanskiĭ, \emph{Proc. Steklov Inst. Math.}, 2000, no. 4 (231), 128--203. No article title or direct link is supplied, and the reference has not been consulted here.
\end{itemize}
\end{document}
